\documentclass[12pt]{article}
\usepackage{amsmath, amssymb, geometry, setspace}
\geometry{margin=1in}
\setstretch{1.2}
\setlength{\parindent}{0pt}


\title{Solutions -- Quiz 3\\ \vspace{10pt}
\large ECON 320 -- Introduction to Mathematical Economics}
\author{Instructor: Muhebullah Karimzada}
\date{October 20, 2025}

\begin{document}
\maketitle


\section*{Part A – Multiple Choice (1 mark each)}

\begin{enumerate}
\item \textbf{(c)} $f'(a)=0$ and $f''(a)<0$ for a local maximum.

\item \textbf{(b)} $f'(x)=0$ and $f''(x)>0$ for a local minimum.

\item \textbf{(c)} $f(x)=x^3-6x^2+9x$. $f'(x)=3x^2-12x+9=3(x-1)(x-3)$, so $x=1,3$. $f''(x)=6x-12$; $f''(1)=-6<0$ (max), $f''(3)=6>0$ (min). Minimum at $x=3$.

\item \textbf{(b)} Profit is maximized where marginal \emph{profit} is zero, equivalently $MR=MC$. Given the answers, ``zero'' is correct.

\item \textbf{(b)} For $y=xe^{-x}$, $y'=e^{-x}(1-x)$, set to zero $\Rightarrow x=1$ (maximum).

\item \textbf{(a)} $f(x)=a-bx^2$ with $b>0$ has $f''(x)=-2b<0$, strictly concave.

\item \textbf{(c)} The necessary condition for an interior extremum is $f'(x)=0$.
\end{enumerate}



\section*{Part B – Analytical Problems (13 marks)}

\subsection*{Question 1 (4 marks)}
A firm’s profit: $\displaystyle \pi(Q)=100Q-4Q^2-20$.

\textbf{(i) FOC.} 
\[
\pi'(Q)=100-8Q.
\]

\textbf{(ii) Maximizer $Q^*$.} Set $\pi'(Q)=0 \Rightarrow 100-8Q=0 \Rightarrow Q^*=\frac{100}{8}=12.5$.

\textbf{(iii) Maximum profit.}
\[
\pi(12.5)=100(12.5)-4(12.5)^2-20=1250-625-20=605.
\]

\textbf{(iv) SOC.} $\pi''(Q)=-8<0$, so $Q^*$ is a strict maximum.

\medskip
\textit{Interpretation.} Output increases profit up to $Q=12.5$; beyond it, diminishing marginal profit turns total profit down.

\bigskip

\subsection*{Question 2 (4 marks)}
Utility: $\displaystyle U(x)=10x-\tfrac12 x^2$.

\textbf{(i) FOC.} $U'(x)=10-x$.

\textbf{(ii) Maximizer $x^*$.} $10-x=0 \Rightarrow x^*=10$.

\textbf{(iii) Value at optimum.} $U(10)=10\cdot10-\tfrac12(10)^2=100-50=50$.

\textbf{(iv) SOC.} $U''(x)=-1<0$, so $x^*$ is a maximum.

\medskip
\textit{Interpretation.} Marginal utility declines linearly; optimal consumption balances benefit and the rising marginal disutility.

\bigskip

\subsection*{Question 3 (5 marks)}
Revenue: $\displaystyle R(P)=500P\,e^{-0.1P}$.

\textbf{(i) FOC.} 
\[
R'(P)=500\left[e^{-0.1P}+P(-0.1)e^{-0.1P}\right]
=500e^{-0.1P}(1-0.1P).
\]
Set $R'(P)=0 \Rightarrow 1-0.1P=0 \Rightarrow P^*=10$.

\textbf{(ii) Maximum revenue.}
\[
R(10)=500\cdot 10 \cdot e^{-1}=\frac{5000}{e}\approx 1839.4.
\]

\textbf{(iii) SOC.} Differentiate again:
\[
R''(P)=500e^{-0.1P}\big(-0.1(1-0.1P)-0.1\big)=500e^{-0.1P}(-0.2+0.01P).
\]
At $P=10$: $R''(10)=500e^{-1}(-0.1)<0$, so a strict maximum.

\textbf{(iv) Economic reasoning.} At low prices, raising $P$ increases revenue; beyond $P=10$, the exponential demand response dominates, quantity falls fast enough that total revenue declines.

\end{document}
